\documentclass[11pt,a4paper]{article}

\usepackage[margin=2.2cm]{geometry}
\usepackage[T1]{fontenc}
\usepackage[utf8]{inputenc}
\usepackage[french]{babel}
\usepackage{lmodern}
\usepackage{microtype}
\usepackage{amsmath,amssymb,amsthm}
\usepackage{booktabs}
\usepackage{xcolor}
\usepackage{enumitem}
\usepackage[most]{tcolorbox}
\usepackage{tikz}
\usetikzlibrary{arrows.meta,positioning}
\usepackage[numbers,sort&compress]{natbib}
\usepackage[colorlinks=true,linkcolor=blue!50!black,citecolor=blue!50!black,urlcolor=blue!50!black]{hyperref}

\setlist{itemsep=0.15em,topsep=0.3em}
\setlength{\parskip}{0.35em}
\setlength{\parindent}{0pt}
\setlength{\emergencystretch}{2em}

\newcommand{\E}{\mathbb{E}}
\newcommand{\code}[1]{\texttt{#1}}
\newcommand{\slide}[2]{\textcolor{gray}{[\,\code{#1}, slide #2\,]}}

\newtcolorbox{codebox}{colback=blue!3,colframe=blue!40!black,boxrule=0.4pt,arc=1pt,
  left=4pt,right=4pt,top=2pt,bottom=2pt,fonttitle=\bfseries\small,title=Dans le code}
\newtcolorbox{keybox}{colback=orange!4,colframe=orange!60!black,boxrule=0.4pt,arc=1pt,
  left=4pt,right=4pt,top=2pt,bottom=2pt,fonttitle=\bfseries\small,title=À retenir}

\title{Mini-projet 1 -- Notes de fond\\\large Biais de surestimation, TD3 et Layer Normalization}
\author{Notes personnelles (non rendues) -- Théo Fontaine}
\date{Octobre 2026}

\begin{document}
\maketitle

\noindent Ces notes supposent le cours lu. Elles ne réexpliquent pas les MDP, Bellman, DQN ou DDPG :
elles renvoient aux slides (dossier \code{RL/}) et approfondissent seulement ce que le projet utilise
au-delà du cours. Les renvois \slide{ddpg.pdf}{n} désignent
\code{2 - DQN, DDPG et TD3/Slides/02 - DDPG and TD3 - slides.pdf}, \code{maxQ.pdf} et
\code{sac\_to\_rlpd.pdf} sont dans \code{5 - Leveraging data/Slides/}.

\tableofcontents

% =====================================================================
\section{Rappel express : DDPG et TD3 dans la notation du cours}

On note $Q_\phi$ pour $\hat Q^{\pi_\theta}_\phi$ (le critique), $\pi_\theta$ l'acteur déterministe,
$\bar\phi$ et $\bar\theta$ les paramètres des réseaux cibles, $d$ l'indicateur « $s'$ est terminal ».

\paragraph{DDPG} \slide{ddpg.pdf}{7, 9, 10}. Critique régressé sur
$y = r + \gamma(1-d)\,Q_{\bar\phi}(s', \pi_{\bar\theta}(s'))$, acteur mis à jour par le gradient
déterministe $\nabla_\theta \E_s[Q_\phi(s,\pi_\theta(s))]$ \citep{silver2014dpg}, cibles mises à jour
en douceur : $\bar\phi \leftarrow \tau\phi + (1-\tau)\bar\phi$.

\paragraph{TD3} \slide{ddpg.pdf}{16, 17}. Trois modifications \citep{fujimoto2018td3} :
\begin{align}
  \tilde a' &= \mathrm{clip}\big(\pi_{\bar\theta}(s') + \mathrm{clip}(\epsilon,-c,c),\,-1,\,1\big),
     \quad \epsilon\sim\mathcal N(0,\sigma^2) && \text{(lissage de la politique cible)}\\
  y &= r + \gamma(1-d)\min_{i=1,2} Q_{\bar\phi_i}(s', \tilde a') && \text{(clipped double-Q)}\\
  L(\phi_i) &= \E\big[(Q_{\phi_i}(s,a) - y)^2\big],\ i=1,2 && \text{(les deux critiques, même cible)}
\end{align}
et l'acteur (qui maximise $Q_{\phi_1}$) ainsi que les trois cibles ne sont mis à jour qu'une fois
tous les \code{policy\_delay} pas de gradient du critique.

\begin{codebox}
\begin{tabular}{@{}ll@{}}
  éq.~1--3, acteur, délai, soft update & \code{1\_implementation.ipynb}, §3 (\code{Agent.update}) \\
  réseau avec ou sans LayerNorm & \code{1\_implementation.ipynb}, §2 (\code{mlp}, \code{Critic}) \\
  retours Monte Carlo, biais & \code{1\_implementation.ipynb}, §4 (\code{evaluate}) \\
  courbes, IC, tests & \code{2\_analysis.ipynb}
\end{tabular}
\end{codebox}

\paragraph{Un seul code pour DDPG et TD3.} Dans le projet, DDPG est \emph{TD3 avec les trois
mécanismes désactivés} : un seul critique, \code{policy\_delay = 1}, pas de bruit sur la politique
cible. On garde un acteur cible dans les deux cas (comme le DDPG original et le « Our DDPG » de
Fujimoto et al.). Ainsi, la seule différence entre les deux algorithmes est exactement ce qu'on
étudie, et tout le reste (réseaux, optimiseur, buffer, évaluation) est identique. Attention : le
notebook du cours utilise $\pi_\theta(s')$ (acteur courant) dans la cible de DDPG ; c'est un choix
légitime mais différent, à mentionner dans le rapport.

% =====================================================================
\section{Le biais de surestimation en profondeur}

\subsection{L'origine : maximiser des estimations bruitées}
Le cours l'affirme \slide{ddpg.pdf}{15} ; voici pourquoi. Supposons que le critique estime sans biais
chaque valeur, $\hat Q(s,a) = Q(s,a) + e_a$ avec $\E[e_a]=0$. Comme $\max$ est une fonction convexe,
l'inégalité de Jensen donne
\begin{equation}
  \E\Big[\max_a \hat Q(s,a)\Big] \;\ge\; \max_a \E\big[\hat Q(s,a)\big] = \max_a Q(s,a).
\end{equation}
Le maximum d'estimations non biaisées est biaisé vers le haut : on sélectionne préférentiellement
les actions dont l'erreur est positive \citep{thrun1993issues}.

\textbf{Exemple chiffré.} Deux actions de vraie valeur $0$, erreurs indépendantes
$e_a\sim\mathcal N(0,\sigma^2)$ : $\E[\max(e_1,e_2)] = \sigma/\sqrt\pi \approx 0{,}56\,\sigma$.
Avec $m$ actions, le biais croît comme $\sigma\sqrt{2\ln m}$ : plus il y a de choix, pire c'est.

\textbf{Et en actions continues ?} DDPG ne calcule pas de $\max$ explicite, mais l'acteur fait une
montée de gradient sur $Q_\phi$ : il cherche $\arg\max_a Q_\phi(s,a)$ \slide{maxQ.pdf}{4, 10}. Il va
donc précisément là où l'erreur du critique est positive. Fujimoto et al. montrent que le biais
existe bien dans ce cadre (leur section 4.1) : c'est le même mécanisme, avec une optimisation
approchée à la place du $\max$.

\subsection{La propagation : le bootstrap transforme une erreur locale en erreur globale}
Le biais ne reste pas local : la cible $y$ contient $Q_{\bar\phi}(s',\cdot)$, donc la surestimation
en $s'$ est régressée dans $Q(s,a)$, puis dans les prédécesseurs de $s$, etc. Si chaque backup ajoute
un biais $\varepsilon$, au point fixe
\begin{equation}
  \hat Q = r + \gamma(\hat Q' + \varepsilon) \quad\Longrightarrow\quad
  \text{biais accumulé} \approx \frac{\varepsilon}{1-\gamma} = 100\,\varepsilon \text{ pour } \gamma=0{,}99.
\end{equation}
S'y ajoute une boucle avec l'acteur : il est attiré par les zones surestimées, on y collecte des
données, mais le critique y généralise mal… Le cours décrit ce \emph{cercle vicieux} pour le RL
offline \slide{sac\_to\_rlpd.pdf}{30} ; en ligne, il est atténué parce qu'on finit par échantillonner
les actions surestimées, ce qui corrige le critique.

\begin{center}
\begin{tikzpicture}[node distance=2.6cm, every node/.style={draw, rounded corners, align=center,
  font=\small, minimum height=0.9cm}, >={Stealth}]
  \node (err) {erreur positive\\de $Q_\phi$ en $(s,a)$};
  \node (act) [right=of err] {l'acteur se\\déplace vers $a$};
  \node (tgt) [right=of act] {$y$ contient la\\surestimation};
  \draw[->] (err) -- (act);
  \draw[->] (act) -- (tgt);
  \draw[->] (tgt) to[bend left=18] node[draw=none, below, font=\footnotesize] {régression du critique : l'erreur grossit} (err);
\end{tikzpicture}
\end{center}

\subsection{Double Q-learning : découpler sélection et évaluation}
L'idée de \citet{vanhasselt2010double} : si l'action est \emph{choisie} avec un estimateur et
\emph{évaluée} avec un autre dont les erreurs sont indépendantes, le biais de sélection disparaît.
DQN $\to$ Double DQN le fait avec le réseau cible. Dans DDPG, l'acteur choisit et le critique cible
évalue, ce qui ressemble à un découplage, mais le critique cible est une moyenne glissante très
lente du critique courant ($\tau=0{,}005$) : leurs erreurs sont presque identiques, donc le
découplage est inefficace. C'est la motivation directe des deux critiques de TD3.

% =====================================================================
\section{TD3 vu par ce projet}

\paragraph{Clipped double-Q : échanger surestimation contre sous-estimation.} Avec deux estimateurs
$Q + e_1$ et $Q + e_2$, prendre le $\min$ biaise vers le bas : avec des erreurs gaussiennes
indépendantes, $\E[\min(e_1,e_2)] = -\sigma/\sqrt\pi$. TD3 assume ce choix : une valeur
\emph{sous-estimée} n'est pas recherchée par l'acteur, donc elle ne se propage pas de la même façon
qu'une surestimation. Le cours signale que la sous-estimation existe bien \slide{ddpg.pdf}{18}, et
RLPD note qu'un seul critique suffit parfois \slide{sac\_to\_rlpd.pdf}{49}. Dans nos mesures, il faut
donc s'attendre à un biais \textbf{proche de zéro, voire négatif} pour TD3.

\paragraph{Lissage de la politique cible : régularisation dans l'espace des actions.} En moyennant
$Q_{\bar\phi}$ sur un petit voisinage de $\pi_{\bar\theta}(s')$, la cible ne peut plus exploiter un
pic étroit d'erreur du critique. Conséquence subtile pour la mesure : le critique de TD3 estime
alors la valeur de la politique \emph{bruitée}, pas exactement celle de $\pi_\theta$. Avec un bruit
de $0{,}2$, l'écart est faible, mais il vaut la peine de le mentionner.

\paragraph{Mises à jour retardées.} L'acteur maximise $Q_{\phi_1}$ : si on le met à jour alors que
le critique est encore très faux, il suit l'erreur. Attendre \code{policy\_delay} pas de critique
réduit la variance de la cible et laisse l'erreur diminuer. Le paramètre est le
\emph{critic-to-actor ratio} de l'énoncé.

\begin{keybox}
Prédiction testable : biais(DDPG) $>$ biais(TD3), surtout tôt dans l'apprentissage. Si tes mesures
ne montrent pas ça, suspecte d'abord un bug (dans la cible ou dans la mesure) avant de conclure.
\end{keybox}

% =====================================================================
\section{Layer Normalization}

\subsection{Définition}
Pour un vecteur de pré-activations $h\in\mathbb R^d$ d'\emph{un} échantillon \citep{ba2016layernorm} :
\begin{equation}
  \mu = \frac1d\sum_{j=1}^d h_j,\qquad \sigma^2 = \frac1d\sum_{j=1}^d (h_j-\mu)^2,\qquad
  \mathrm{LN}(h) = \gamma\odot\frac{h-\mu}{\sqrt{\sigma^2+\epsilon}} + \beta,
\end{equation}
avec $\gamma,\beta\in\mathbb R^d$ appris (initialisés à $1$ et $0$). En PyTorch :
\code{nn.LayerNorm(d)}.

\paragraph{LN vs BatchNorm.} BatchNorm \citep{ioffe2015batchnorm} normalise chaque neurone sur le
\emph{batch} : la sortie pour un échantillon dépend des autres, et le comportement diffère entre
entraînement (statistiques du batch) et inférence (moyennes glissantes). En RL off-policy c'est
problématique : on évalue $Q(s,a)$ sur les données du buffer, $Q(s',\tilde a')$ sur des actions qui
n'y sont pas, et le critique cible est une copie avec d'autres statistiques. CrossQ
\citep{bhatt2024crossq} y arrive, mais avec des précautions spécifiques. LayerNorm normalise
\emph{chaque échantillon indépendamment}, n'a pas d'état, et se comporte pareil partout : il
s'insère sans précaution dans un critique. Le cours évoque d'ailleurs des résultats non concluants de
BatchNorm dans DDPG \slide{ddpg.pdf}{5}.

\subsection{Pourquoi LN limite la surestimation}
Le cours le présente comme une recette de RLPD \slide{sac\_to\_rlpd.pdf}{40, 45--47} : « sans
LayerNorm, les Q-values sont surestimées ». L'argument de \citet{ball2023rlpd} est une
\textbf{borne}. Notons $\varphi(s,a) = \mathrm{ReLU}(\mathrm{LN}(z))$ la dernière couche cachée et
$Q = w^\top\varphi + b$. Le vecteur normalisé $(z-\mu)/\sigma$ a toujours une norme $\sqrt d$, et
ReLU ne fait que diminuer la norme, donc
\begin{equation}
  |Q_\phi(s,a)| \;\le\; \|w\|_2\,\big(\|\gamma\|_\infty\sqrt d + \|\beta\|_2\big) + |b|
  \qquad\text{pour \emph{toute} entrée }(s,a).
\end{equation}
La valeur maximale du critique est contrôlée par les poids, pas par l'entrée. À l'inverse, un MLP
ReLU sans normalisation est linéaire par morceaux : loin des données, il \emph{extrapole
linéairement}, et peut produire des valeurs arbitrairement grandes que l'acteur s'empressera
d'exploiter. LN empêche cette extrapolation catastrophique hors distribution
\slide{sac\_to\_rlpd.pdf}{46}.

\paragraph{Les autres effets connus.}
\begin{itemize}
  \item \textbf{Plasticité} : les réseaux entraînés sur des cibles non stationnaires perdent
    progressivement leur capacité à apprendre. LN fait partie des remèdes les plus efficaces
    \citep{lyle2024plasticity}.
  \item \textbf{Passage à l'échelle} : BRO \citep{nauman2024bro} et SimBa \citep{lee2025simba}
    utilisent LN pour entraîner de gros critiques et des UTD élevés.
  \item \textbf{UTD élevé} : DroQ \citep{hiraoka2022droq} combine dropout et LayerNorm
    \slide{sac\_to\_rlpd.pdf}{11}.
\end{itemize}

\paragraph{Ce qu'on peut s'attendre à voir sur LunarLander.} Les actions sont bornées par le
$\tanh$ de l'acteur, donc l'extrapolation hors distribution dans l'espace des actions est limitée :
l'effet de LN peut être plus faible que dans les benchmarks de RLPD. C'est précisément une question
empirique, et un résultat « LN ne change pas grand-chose ici » est un résultat valable s'il est
statistiquement établi.

\begin{codebox}
LN est inséré dans chaque bloc caché du \textbf{critique seulement} :
\code{Linear} $\to$ \code{LayerNorm} $\to$ \code{ReLU}, rien après la couche de sortie.
Switch : \code{layernorm} de \code{Config} ; implémenté dans \code{mlp}
(\code{1\_implementation.ipynb}, §2).
\end{codebox}

% =====================================================================
\section{Mesurer le biais}

\subsection{Le principe}
On compare la valeur prédite à la valeur réelle de la politique. On fait tourner la politique
\emph{déterministe} $\pi_\theta$ (sans bruit d'exploration) et, pour chaque pas $t$ d'un épisode de
longueur $T$, on calcule le retour Monte Carlo actualisé
\begin{equation}
  G_t = \sum_{k=0}^{T-t-1}\gamma^k\, r_{t+k+1}, \qquad\text{récursivement : } G_{T-1} = r_T,\
  G_t = r_{t+1} + \gamma\, G_{t+1},
\end{equation}
qui est un estimateur sans biais de $Q^{\pi}(s_t,a_t)$. On le compare à $Q_{\phi_1}(s_t,a_t)$ (le
critique que l'acteur maximise).

\textbf{Piège classique :} $Q$ est une quantité \emph{actualisée}. On la compare à $G_t$, pas au
retour d'épisode non actualisé affiché par l'évaluateur (de l'ordre de $+200$ pour un atterrissage
réussi, alors que $G_0$ avec $\gamma=0{,}99$ est bien plus petit).

\subsection{Normalisation}
Comme l'échelle des retours change au cours de l'apprentissage, on suit \citet{chen2021redq} :
\begin{equation}
  \text{biais normalisé} = \frac{Q_{\phi_1}(s_t,a_t) - G_t}{\big|\,\E_{t}[G_t]\,\big|}.
\end{equation}
On en rapporte la \textbf{moyenne} (le biais : surestimation si $>0$) et l'\textbf{écart-type}
(la dispersion de l'erreur, qui pénalise aussi la politique). On garde aussi le biais brut
$\E[Q - G]$, plus facile à interpréter physiquement.

\subsection{Troncature}
LunarLander tronque les épisodes à $1000$ pas. Si un épisode est \emph{tronqué} (et non terminé), le
$G_t$ calculé ignore tout ce qui se passe après $T$, alors que $Q$ en tient compte : les derniers
$G_t$ sont faux. La queue manquante pèse $\gamma^{T-t}$. On ne garde donc que les pas où elle est
négligeable :
\begin{equation}
  t < T - H \quad\text{avec}\quad \gamma^H \le \text{tol}
  \;\Longleftrightarrow\; H = \Big\lceil \frac{\ln \text{tol}}{\ln\gamma}\Big\rceil
  \qquad (\gamma=0{,}99,\ \text{tol}=0{,}01 \Rightarrow H = 459).
\end{equation}
Si l'épisode est \emph{terminé} (crash ou atterrissage), la valeur après $T$ est réellement nulle et
tous les pas sont valides. C'est exactement la distinction \code{terminated}/\code{truncated} du
notebook du cours, appliquée à la mesure.

\subsection{Limites à discuter dans le rapport}
\begin{itemize}
  \item On mesure le biais sur la distribution d'états de $\pi_\theta$, depuis les états initiaux,
    et non sur des états tirés du buffer (cela demanderait de réinitialiser l'environnement dans un
    état arbitraire).
  \item $G_t$ est un estimateur bruité (états initiaux et terrain aléatoires) : on moyenne sur
    plusieurs épisodes et sur les seeds.
\end{itemize}

\begin{codebox}
\code{1\_implementation.ipynb}, §4 : \code{discounted\_returns} (récursion ci-dessus),
\code{valid\_length} (troncature), et \code{evaluate}, qui fait les rollouts et calcule le biais
normalisé.
\end{codebox}

% =====================================================================
\section{Les autres hyper-paramètres mentionnés par l'énoncé}

\begin{description}
  \item[UTD ratio (update-to-data).] Nombre de pas de gradient par transition collectée
    \slide{sac\_to\_rlpd.pdf}{9--10, 48}. Plus d'updates = meilleure efficacité en échantillons, mais
    plus de surapprentissage du buffer et plus de surestimation \citep{li2023overfitting} : chaque
    pas de gradient refait un backup biaisé. Prédiction : le biais augmente avec l'UTD, et LN
    devrait aider davantage à UTD élevé.
  \item[Critic-to-actor ratio.] C'est \code{policy\_delay} de TD3 (section 3).
  \item[Profondeur (2 ou 3 couches cachées).] Plus de capacité, mais des réseaux plus profonds non
    normalisés sont plus instables ; c'est le cas où LN est censé le plus aider (BRO, SimBa).
\end{description}
Conseil de l'énoncé : « a small accurate study is better than a large inaccurate one ». On fait
d'abord la grille principale (DDPG/TD3 $\times$ LN) proprement, et un facteur supplémentaire
seulement si le temps de calcul le permet.

% =====================================================================
\section{Méthodologie expérimentale}

\paragraph{Pourquoi plusieurs seeds.} En deep RL, deux runs d'un même algorithme avec des seeds
différentes peuvent différer plus que deux algorithmes différents \citep{henderson2018matters}. Une
courbe sur une seed ne prouve rien. Les seeds fixent tout : initialisation des réseaux, bruit
d'exploration, environnements. On rapporte toujours une statistique \emph{sur les seeds} : chaque
seed est \emph{un} échantillon, et les épisodes d'évaluation d'une même seed n'en sont pas des
indépendants.

\paragraph{Combien de seeds.} \citet{colas2018seeds} (dont l'auteur est notre enseignant) étudient la
puissance statistique : avec 5 seeds, on ne détecte de façon fiable que des différences
importantes. 5 est un minimum, 10 c'est mieux si le calcul le permet. Il faut dire dans le rapport
qu'un test non significatif avec peu de seeds ne prouve pas l'absence d'effet.

\paragraph{Intervalles de confiance par bootstrap.} Pour une configuration et un instant donnés, on a
$n$ valeurs (une par seed). On les rééchantillonne avec remise $B$ fois (ex. $B=10\,000$), on calcule
la moyenne de chaque rééchantillon, et l'IC à 95\,\% est donné par les percentiles 2,5 et 97,5 de ces
moyennes. Aucune hypothèse de normalité. C'est ce que tracera \code{scripts/analyze.py} autour des
courbes moyennes.

\paragraph{Tests.} Pour comparer deux configurations sur une métrique finale (retour, biais) :
\begin{itemize}
  \item \textbf{t-test de Welch} : compare les moyennes sans supposer des variances égales.
    Il suppose une distribution à peu près normale, ce qui est discutable avec 5 points.
  \item \textbf{Mann--Whitney U} : non paramétrique, basé sur les rangs, plus robuste aux valeurs
    aberrantes (un run qui diverge).
\end{itemize}
\citet{colas2019hitchhiker} discutent ces choix en détail. On rapporte la $p$-value \emph{et} la
taille d'effet (différence des moyennes avec son IC), et on fixe $\alpha=0{,}05$ \emph{avant} de
regarder les résultats. Avec plusieurs comparaisons, certaines sortiront significatives par hasard
(correction de Bonferroni : $\alpha/k$ pour $k$ tests). \citet{agarwal2021rliable} proposent des
agrégats robustes (moyenne interquartile, IQM), utiles quand les seeds sont très dispersées.

\paragraph{Le lien biais $\leftrightarrow$ performance (Q2).} Pour chaque run, on trace le biais
moyen contre le retour final, une couleur par configuration, et on calcule une corrélation (Pearson
ou Spearman) avec sa $p$-value. Attention : corrélation n'est pas causalité. Un agent qui ne sait pas
atterrir a aussi des retours faibles qui changent l'échelle du biais.

\paragraph{La figure checklist du cours} (\code{0 - Généralités/01 - The figure checklist.pdf}). Chaque
figure doit avoir :
\begin{itemize}
  \item un titre et une légende autonome, qui dit ce qu'il faut retenir ;
  \item des axes avec labels et échelles ;
  \item si c'est une moyenne, sa variabilité, avec sa nature précisée dans la légende (ici : IC
    bootstrap à 95\,\% sur $n$ seeds) ;
  \item le code couleur s'il y en a un ;
  \item les résultats des tests statistiques dans la figure ou la légende ;
  \item un numéro, et une mention dans le texte.
\end{itemize}

% =====================================================================
\section{Glossaire}
\begin{description}[style=nextline, leftmargin=1.5em]
  \item[Bootstrap (en TD)] Utiliser sa propre estimation $Q(s',\cdot)$ dans la cible. Rien à voir
    avec le bootstrap statistique de la section 7, homonyme malheureux.
  \item[Biais de surestimation] $\E[Q_\phi(s,a)] - Q^\pi(s,a) > 0$, mesuré ici par $Q - G$.
  \item[Clipped double-Q] Cible calculée avec le minimum de deux critiques cibles.
  \item[Hors distribution (OOD)] Paire $(s,a)$ éloignée des données d'entraînement du critique.
  \item[Plasticité] Capacité d'un réseau à continuer d'apprendre de nouvelles cibles.
  \item[Terminated / truncated] Fin réelle de l'épisode (valeur future nulle) / coupure par limite de
    temps (on bootstrape ; dans la mesure MC, la queue manque).
  \item[UTD ratio] Pas de gradient par pas d'environnement.
\end{description}

\bibliographystyle{plainnat}
\bibliography{refs}

\end{document}
